El análisis de algoritmos suele apoyarse en la notación asintótica para
comparar su eficiencia teórica, pero esta no siempre predice su
comportamiento real sobre hardware concreto. En este trabajo se estudian
experimentalmente dos problemas clásicos: el ordenamiento de arreglos
unidimensionales de números enteros, mediante los algoritmos
\textsc{MergeSort}, \textsc{QuickSort}, \textsc{PatienceSort} y
\texttt{std::sort} y la multiplicación de matrices cuadradas,
mediante el algoritmo \textsc{Naive} ($O(n^3)$) y el algoritmo de
\textsc{Strassen} ($O(n^{\log_2 7})$). El objetivo es medir tiempo de
ejecución y uso de memoria de cada algoritmo bajo distintos tamaños de
entrada, tipos de dato (ordenado, invertido, aleatorio; denso, diagonal,
disperso) y dominios de valores, para luego contrastar estos resultados
con su complejidad teórica y entender en qué medida esta última describe
o no su desempeño práctico.